\documentclass{article}


\usepackage[preprint]{utils_packages/neurips_2023}



\usepackage[utf8]{inputenc} % allow utf-8 input
\usepackage[T1]{fontenc}    % use 8-bit T1 fonts
\usepackage{hyperref}       % hyperlinks
\usepackage{url}            % simple URL typesetting
\usepackage{booktabs}       % professional-quality tables
\usepackage{amsfonts}       % blackboard math symbols
\usepackage{nicefrac}       % compact symbols for 1/2, etc.
\usepackage{microtype}      % microtypography
\usepackage{xcolor}         % colors
\usepackage{float}
\usepackage{listings}

\urlstyle{same}
\hypersetup{
  hidelinks,
  pdfcreator={LaTeX via pandoc}}

\author{}
\date{}

% [Custom packages] from here down
% [Custom packages] from here down 
\usepackage{float}
\usepackage{todonotes}
\usepackage{amsmath}
\usepackage{algorithm}
\usepackage{algpseudocode}
\usepackage{enumitem}
\usepackage{caption}
\usepackage{subcaption}
% \usepackage{biblatex}
% \addbibresource{refs.bib}

\newcommand{\x}{\boldsymbol{x}}
% \newcommand{\y}{\boldsymbol{y}}
\newcommand{\X}{\boldsymbol{X}}
\newcommand{\w}{\boldsymbol{w}}
\newcommand{\sign}{\text{sign}}
\newcommand{\ub}{\boldsymbol{u}}
\newcommand{\z}{\boldsymbol{z}}
\newcommand{\vb}{\boldsymbol{v}}

\newcommand{\numberset}{\mathbb}
\newcommand{\N}{\numberset{N}}
\newcommand{\Z}{\numberset{Z}}
\newcommand{\Q}{\numberset{Q}}
\newcommand{\R}{\numberset{R}}
\newcommand{\C}{\numberset{C}}
\newcommand{\E}{\numberset{E}}
\newcommand{\Prob}{\numberset{P}} 

\usepackage{mathtools}
\DeclarePairedDelimiter\ceil{\lceil}{\rceil}
\DeclarePairedDelimiter\floor{\lfloor}{\rfloor}
\newcommand{\norm}[1]{\left\lVert#1\right\rVert}

\usepackage{nomencl}
\def\bmat#1{\left[\begin{array}{#1}}
\def\emat{\end{array}\right]}
\def\bv {\bmat{c} }
\def\ev {\emat}

\newcommand {\bsis} {\left\{ \begin{array} }
\newcommand {\esis} {\end{array}\right.}

\def\Sum#1#2{\sum\limits_{#1}^{#2}}
\def\set#1#2{\{ \; #1 \;:\;#2\;\}} % creates a set: { #1 : #2 }



\newcommand{\y}{\hat{y}}

\newcommand{\tauB}{\boldsymbol{\tau}}
\newcommand{\xiB}{\boldsymbol{\xi}}
\newcommand{\thetaB}{\boldsymbol{\theta}}
\newcommand{\thB}{{\boldsymbol{\theta}}}
\newcommand{\alB}{{\boldsymbol{\alpha}}}
\newcommand{\betB}{{\boldsymbol{\beta}}}
\newcommand{\yB}{{\boldsymbol{y}}}


\setcounter{MaxMatrixCols}{20}
\newtheorem{remark}{Remark}[section]


% [Custom packages] from here up 

\usepackage{amsmath}

\usepackage{comment}

% [Custom packages] from here up


\title{EvoScale: Distributed Evolutionary Search\\
  with Online Reinforcement Learning\\[0.3em]
  \large Midterm Check-in}

\author{
Group 7 \\ Matteo Guarrera \quad Elizabeth Polito \quad Harper Lyu \\ Mert Cemri \quad Sultan Daniels \quad Eric Xu
}

\begin{document}
\maketitle
\author{}
\vspace{-0.5cm}

% -----------------------------------------------------------------------
\section*{Project Summary}

EvoScale extends AdaEvolve~\cite{adaevolve2025}, the current state of the art in LLM-driven evolutionary program optimization, along two axes: (1)~massively parallel multi-problem execution with KV cache sharing to maximize GPU utilization, and (2)~online reinforcement learning (GRPO) to fine-tune the LLM mutation operator during search. Our central thesis is that at sufficient scale, inference is training: every token generated during search also produces a gradient signal that improves subsequent search iterations.

At the proposal stage, AdaEvolve had been evaluated exclusively with proprietary frontier models (GPT-5, Gemini-3-Pro). A key prerequisite for EvoScale---which requires weight access for RL fine-tuning---was demonstrating that AdaEvolve's adaptive mechanisms work with open-source models. This report presents that milestone and outlines our path forward.

% -----------------------------------------------------------------------
\section*{Preliminary Results}

\subsection*{Open-Source Backbone Validation}

Since the proposal, we have completed a set of experiments validating AdaEvolve with open-source models. We ran two Qwen3.5 variants (9B and 27B parameters) served locally via vLLM on a single H100 GPU, using AdaEvolve's standard configuration with \emph{zero hyperparameter changes} from the frontier-model experiments. Table~\ref{tab:backbone} summarizes results across four benchmarks (100 iterations, 3 seeds each).

\begin{table}[h]
  \centering
  \small
  \setlength{\tabcolsep}{3.5pt}
  \renewcommand{\arraystretch}{1.15}
  \begin{tabular}{l cc cc cc cc}
  \toprule
  & \multicolumn{2}{c}{\textbf{Circle Packing}}
  & \multicolumn{2}{c}{\textbf{Signal Proc.}}
  & \multicolumn{2}{c}{\textbf{PRISM}}
  & \multicolumn{2}{c}{\textbf{TXN Sched.}} \\
  \cmidrule(lr){2-3}
  \cmidrule(lr){4-5}
  \cmidrule(lr){6-7}
  \cmidrule(lr){8-9}
  \textbf{Backbone}
  & Avg & Best
  & Avg & Best
  & Avg & Best
  & Avg & Best \\
  \midrule
  GPT-5        & \textbf{2.629} & \textbf{2.636} & \textbf{0.698} & \textbf{0.718} & \textbf{26.30} & \textbf{26.37} & \textbf{4,317} & \textbf{4,348} \\
  Gemini-3-Pro & 2.632 & \textbf{2.636} & 0.593 & 0.603 & 26.26 & 26.26 & 4,221 & 4,310 \\
  Qwen3.5-27B  & 2.626 & 2.630 & 0.662 & 0.695 & 25.13 & 25.23 & 3,905 & 4,082 \\
  Qwen3.5-9B   & 2.253 & 2.607 & 0.634 & 0.665 & 24.92 & 25.88 & 3,864 & 4,098 \\
  \bottomrule
  \end{tabular}
  \caption{AdaEvolve performance across LLM backbones (Mean and Best over 3 seeds, 100 iterations). Higher is better for all metrics.}
  \label{tab:backbone}
\end{table}

Three findings stand out:

\textbf{1. The 27B model is surprisingly competitive.} On Circle Packing, Qwen3.5-27B reaches 2.626 average (vs.\ GPT-5's 2.629)---a gap of only 0.1\%. On the best run, it reaches 2.630, within 0.2\% of the best-known result. This is remarkable for a model roughly 10$\times$ smaller than frontier APIs and running on a single GPU.

\textbf{2. AdaEvolve's adaptive control transfers without modification.} We used identical hyperparameters ($\rho=0.9$, $\tau_S=0.02$, $\tau_M=0.12$, etc.) across all four backbones. The fact that the same adaptive thresholds work for both a 9B open-source model and GPT-5 confirms that the improvement signal $G_t$ is backbone-agnostic---it tracks the search dynamics, not the model's absolute capability.

\textbf{3. The 9B model shows high variance but non-trivial performance.} The 9B variant achieves 2.607 best-of-3 on Circle Packing (within 1.1\% of GPT-5's best) despite an average of 2.253 with std of 0.365. This high variance is precisely the kind of signal that online RL should help stabilize: the model \emph{can} produce strong mutations but does so inconsistently, suggesting that fine-tuning on successful trajectories could substantially improve its reliability.

\subsection*{What These Results Mean for EvoScale}

These experiments de-risk the core assumption of the project: that open-source models are viable backbones for evolutionary search. The Qwen3.5-27B model is strong enough to serve as the GRPO base policy (starting from near-frontier performance rather than bootstrapping from zero), the 9B model's high variance makes it an ideal testbed for whether RL can stabilize performance, and vLLM serving on a single H100 validates the infrastructure path for live weight updates.

We note that the backbone validation benchmarks (Circle Packing, Signal Processing, PRISM, TXN) differ from the proposed evaluation benchmarks (ARC-AGI-2, Triton, AtCoder). This was intentional: these four tasks are AdaEvolve's standard suite with established baselines, making them the fastest way to validate the backbone. The proposed evaluation benchmarks remain the target for the full EvoScale system; ARC-AGI-2 evaluation harnesses are set up and we have confirmed the 120-task evaluation split runs end-to-end with the 27B model.

\subsection*{Implementation Progress}

Beyond the backbone experiments, we have made progress on the systems infrastructure:
\begin{itemize}
    \item \textbf{Ray orchestration (prototype):} We have a working prototype that launches $N$ independent AdaEvolve actors sharing a single vLLM endpoint. Currently tested with $N=10$ problems on a single H100; scaling to $N=50{-}200$ is the next step.
    \item \textbf{Prefix caching verified:} We confirmed that vLLM's radix attention cache achieves $>$95\% hit rate on repeated system prompts across actors, validating the KV cache sharing assumption (details in Methodology Refinement below).
    \item \textbf{GRPO loop:} Not yet started. This is the primary focus for the next two weeks.
\end{itemize}

From the existing logs, we measured the mutation improvement rate across the four benchmarks with the 27B model: 12.3\% of mutations yield positive fitness improvement (averaged across tasks and seeds), consistent with our proposal estimate of 10--15\%. With $K=8$ candidates per iteration, this gives approximately 1 positive signal per iteration, which is sufficient for GRPO but leaves limited margin. We plan to add parse-correctness as a dense auxiliary reward to supplement the sparse improvement signal.

% -----------------------------------------------------------------------
\section*{Methodology Refinement}

Based on the open-source experiments, we have refined the project plan in two ways:

\textbf{KV cache sharing is even more important than anticipated.} When serving Qwen3.5-27B on a single H100, the KV cache for a 10,000-token ARC-AGI-2 system prompt consumes approximately 2.5~GB of GPU memory. With $K=8$ parallel candidates per iteration, naive serving would require 20~GB of KV cache for identical prefixes. vLLM's radix attention prefix cache eliminates this redundancy entirely: one copy of the system prompt KV cache is shared across all 8 candidates, reducing memory from 20~GB to 2.5~GB. This is not merely an optimization---it is a \emph{prerequisite} for fitting the 27B model with batch generation on a single GPU. Without prefix caching, we would need to either reduce batch size (hurting GRPO signal quality) or use multiple GPUs (increasing cost). We plan to measure and report the exact memory savings and throughput improvement from prefix caching as part of our systems evaluation.

\textbf{We are starting with the 27B model, not 32B.} The proposal targeted Qwen2.5-Coder-32B, but the 27B variant from the newer Qwen3.5 family achieves better search performance in our experiments (likely due to improved instruction following). We will use Qwen3.5-27B as the base policy for GRPO, with Qwen3.5-9B as a secondary testbed for studying RL's effect on high-variance models.

% -----------------------------------------------------------------------
\section*{Risk Assessment}

We identified two key risks at proposal time. Here is how our preliminary results inform each:

\textbf{Risk 1: Open-source models may be too weak for evolutionary search.} \emph{Status: Mitigated.} The 27B model reaches 98--99\% of GPT-5 performance on Circle Packing and Signal Processing. The gap is larger on TXN Scheduling (3,905 vs.\ 4,317, roughly 90\%), but the model still produces meaningful solutions that could serve as a strong starting point for RL fine-tuning. The 9B model's high variance is a concern for direct deployment but is actually \emph{favorable} for RL: high variance means the policy explores a wide distribution of mutations, giving GRPO diverse comparison signals within each batch.

\textbf{Risk 2: Reward sparsity may prevent effective RL training.} \emph{Status: Partially mitigated.} With the 27B model achieving non-trivial scores on all four benchmarks, we expect 10--15\% of mutations to yield positive improvements, consistent with our proposal estimates. However, we have not yet measured the actual improvement rate from the open-source model runs. This is a key metric to extract from the existing logs before starting RL experiments.

\textbf{New risk: Memory pressure from concurrent LoRA training and serving.} Running LoRA fine-tuning while serving inference on the same GPU introduces memory contention. With the 27B model consuming approximately 54~GB in FP16 for weights alone, plus KV cache and LoRA adapter gradients, a single 80~GB H100 may be tight. Our mitigation plan is to use 4-bit quantization for the base model weights during serving (reducing footprint to $\sim$14~GB) while keeping the LoRA adapter in FP16 for gradient quality. vLLM's weight-update API supports this mixed-precision configuration.

% -----------------------------------------------------------------------
\section*{Potential Impact}

The open-source model results strengthen our original impact claims:

\textbf{Closing the OSS-to-frontier gap.} Qwen3.5-27B already reaches 98--99\% of GPT-5 on 2 of 4 benchmarks without any fine-tuning. The remaining experiments will test whether online GRPO can close the gap on harder tasks (TXN: currently 90\% of GPT-5) or exceed frontier performance on specific benchmarks.

\textbf{KV cache optimization as a scaling enabler.} Our experiments confirm that prefix caching is essential, not optional, for parallel evolutionary search with open-source models. The 20~GB $\rightarrow$ 2.5~GB reduction for $K=8$ parallel candidates is the difference between fitting on one GPU and needing a multi-GPU setup.

\textbf{Variance reduction as a testable hypothesis.} The 9B model's high variance (std 0.365 on Circle Packing) alongside occasional strong performance (best-of-3: 2.607) provides a concrete, measurable test for whether RL training stabilizes evolutionary search. If GRPO reduces variance while maintaining peak performance, the improvement will be clearly visible in the numbers --- not a speculative claim about scaling laws but a direct before/after comparison on the same benchmarks.

% -----------------------------------------------------------------------
\section*{Timeline to Final Report (4 weeks remaining)}

\begin{itemize}
    \item \textbf{Week 1:} Scale Ray orchestration from $N=10$ to $N=50$ problems. Run parallel AdaEvolve (no RL) on ARC-AGI-2 with the 27B model. Report throughput, KV cache hit rates, and GPU utilization. This establishes the ``parallel-only'' baseline. \emph{(Matteo, Eric)}
    \item \textbf{Week 2:} Implement GRPO training loop with live LoRA weight updates. First RL experiments on ARC-AGI-2 with the 27B model. \emph{(Mert, Sultan, Harper)}
    \item \textbf{Week 3:} Run ablations (parallel-only vs.\ parallel+RL, 27B vs.\ 9B). If time permits, generalization test: evaluate ARC-AGI-2-trained adapter on Triton kernels. \emph{(Elizabeth, Harper)}
    \item \textbf{Week 4:} Final report, systems benchmarking, analysis. Buffer for debugging RL instabilities from week 2.
\end{itemize}

\textbf{Scoping note:} If GRPO proves unstable within the timeline, we scope the final report to the parallel-only system (Ray + KV cache + multi-problem orchestration) with thorough throughput analysis and ablations, which is a solid systems contribution on its own. The RL component is the stretch goal; the parallel infrastructure is the guaranteed deliverable.

% -----------------------------------------------------------------------
\bibliographystyle{plain}
\begin{thebibliography}{9}

\bibitem{adaevolve2025}
AdaEvolve: Adaptive LLM-Driven Zeroth-Order Optimization.
\textit{arXiv:2602.20133}, 2025.
\url{https://arxiv.org/abs/2602.20133}

\bibitem{deepseekr1}
DeepSeek-AI.
DeepSeek-R1: Incentivizing Reasoning Capability in LLMs via Reinforcement
Learning.
\textit{arXiv:2501.12948}, 2025.

\bibitem{vllm}
W.~Kwon et al.
Efficient Memory Management for Large Language Model Serving with
PagedAttention.
\textit{SOSP}, 2023.

\bibitem{qwen25coder}
Hui et al.
Qwen2.5-Coder Technical Report.
\textit{arXiv:2409.12186}, 2024.

\bibitem{lora}
E.~J.~Hu et al.
LoRA: Low-Rank Adaptation of Large Language Models.
\textit{ICLR}, 2022.

\end{thebibliography}


\end{document}
